% One-time imported draft. Editable native source; no Markdown regeneration.
\section{A calculation as machine
  state}\label{001-console-calculation-p00-a-calculation-as-machine-state}

The expression \texttt{2\ \ensuremath{\times}\ x\ +\ 1} specifies a result. To execute it,
a machine also needs a representation for \texttt{x}, a place for
intermediate values, and instructions that describe each change. The
first program makes these choices visible. It reads an integer into a
register, computes the expression, and passes the result to a console
service.

A register holds a bit pattern. The instruction using that register
determines how those bits participate in an operation. During this
lesson, the useful state is small enough to record in a table: the
input, the doubled value, and the final answer. The debugger lets us
compare those predictions with the running program.

The task accepts inputs from \ensuremath{-}1000 through 1000. Stating that range lets
us check that the demonstrated intermediate values fit in a signed
32-bit word. The program is responsible for the specified calculation
within that range.
